Revenue function:  R(x) = price \times quantity

If price increases by  \Delta p, quantity decreases by  m \Delta p:
Q = Q_0 - m(p - p_0)
R = p \cdot Q = -mp^2 + (Q_0 + mp_0)p

Maximize: vertex at  p^* = \frac{Q_0 + mp_0}{2m}

Example (2024 \#10): 360 dinners at \$50; every \$5 increase \to 20 fewer dinners
p = 50 + 5t, \quad Q = 360 - 20t
R = (50+5t)(360-20t) = -100t^2 + 800t + 18000
Vertex:  t = \frac{800}{200} = 4 \rightarrow p = \$70
R(4) = 70 \times 280 = 19600

Limousine example (2023 \#15): \$10/ride, 300 riders/day
Every \$1 increase \to 15 fewer riders
R = (10+x)(300-15x) = -15x^2 + 150x + 3000
Vertex:  x = \frac{150}{30} = 5 \rightarrow p = \$15

AM-GM for maximizing product given sum:
xy \le \left(\frac{x+y}{2}\right)^2  with equality when  x = y

When  x + y = c,  xy  is maximized when  x = y = \frac{c}{2}:
(xy)_{\max} = \frac{c^2}{4}

Example (2018 \#9):  4x + 5y = 20, maximize  xy
Max at  4x = 5y and 4x+5y=20 \rightarrow x=2.5,\; y=2
xy_{\max} = 5
